\begin{abstract}
Recent progress in pre-trained neural language models has significantly improved the performance of many natural language processing (NLP) tasks. In this paper we propose a new model architecture \textbf{DeBERTa}
%, which stands for 
(\textbf{D}ecoding-\textbf{e}nhanced \textbf{BERT} with disentangled \textbf{a}ttention) 
%Different from the BERT\cite{devlin2018bert} or RoBERTa\cite{liu2019roberta}, 
that improves the BERT and RoBERTa models using two novel techniques.
% the language model pre-training of a masked language model. 
%First, instead of using absolute position embedding to model sequence dependency, DeBERTa applies a new relative position encoding method called disentangled attention which thoroughly captures the dependency between positions and content in the sequence. 
The first is the disentangled attention mechanism, where each word is represented using two vectors that encode its content and position, respectively, and the attention weights among words are computed using disentangled matrices on their contents and relative positions, respectively.
Second, 
%inspired by the encoder-decoder attention in original transformer structure designed for seq-to-seq tasks, 
an enhanced mask decoder is used to incorporate absolute positions in the decoding layer to predict the masked tokens in model pre-training.
%an enhanced mask decoder is used to replace the output softmax layer 
%the last layer of Transformer $decoder$ 
%to decode masked tokens for model pre-training. 
% which encourage the self-attention to capture global contextual information instead of looking at itself.
%on top of existing multi-layer transformer encoder, DeBERTa introduces the notion of multi-layer decoding  to better reconstruct masked tokens. 
In addition, a new virtual adversarial training method is used for fine-tuning to improve models' generalization.  
We show that these techniques significantly improve the efficiency of model pre-training and the performance of both natural language understand (NLU) and natural langauge generation (NLG) downstream tasks. 
Compared to RoBERTa-Large, a DeBERTa model 
% with 45\% model parameters
trained on half of the training data
%computational capacity 
% of the amount used to train the RoBERTa-Large model, our DeBERTa-Large model 
performs consistently better on a wide range of NLP tasks, 
% e.g., 
achieving improvements on MNLI by +0.9\% (90.2\% vs. 91.1\%), on SQuAD v2.0 by +2.3\% (88.4\% vs. 90.7\%) and RACE by +3.6\% (83.2\% vs. 86.8\%). 
Notably, we scale up DeBERTa by training a larger version that consists of 48 Transform layers with 1.5 billion parameters. 
The significant performance boost makes the single DeBERTa model surpass the human performance on the SuperGLUE benchmark~\citep{wang2019superglue} for the first time in terms of macro-average score (89.9 versus 89.8), and the ensemble DeBERTa model sits atop the SuperGLUE leaderboard as of January 6, 2021, outperforming the human baseline by a decent margin (90.3 versus 89.8). The pre-trained DeBERTa models and the source code were released at: \url{https://github.com/microsoft/DeBERTa}\footnote{Our code and models are also available at HuggingFace Transformers: \url{https://github.com/huggingface/transformers}, \url{https://huggingface.co/models?filter=deberta} }.
%We will release the DeBERTa models and the source code to the public.

%The substantially outperforms Google's T5 with 11 billion parameters on the SuperGLUE benchmark \citep{wang2019superglue} and, for the first time, surpasses the human performance (89.9 vs. 89.8). 

% The relative contributions of the two techniques are validated by an ablation study.
%Meanwhile, extensive ablation study is presented to validate the two contributions introduced by DeBERTa.
%Moreover, we continue to run a series of elaborate ablation studies to clearly demonstrate the merits of the two innovations introduced by DeBERTa.  
%The DeBERTa code and pre-trained models will be made publicly available. 

%Recent progress on this area such as BERT \cite{devlin2018bert}, XLNet \cite{yang2019xlnet}, RoBERTa \cite{liu2019roberta} push the pre-training to a new stage with the successfully applying of transformer and masked language model. However, they usally requires a lot of training data, parameters and a lot of training steps which means there are huge space to improve their efficiency. In this work, we explored several ways to improve the efficiency of the per-training of MLM with significant reduced training time and data requirement while the result is much better than previous results.
\end{abstract}

